% Immutable cited-input supplement. Target proofs remain embargoed.
% Target: night:P4.9; tag 02Y1.
% Frozen source SHA256 74497bfe327fae25c9cdcb17334b3bf42a9793a9df7c0a4404d97568f2f760f0.
% CompendiumRef is resolved against the cited S6_09 source edition, not silently identified with the current series statement.

% Exact cited input: series:II / night:P4.26
\begin{proposition}[P4.26]\hypertarget{tag-02WY}{}\label{night:P4.26}\label{stmt:P4.26}


In Part~4.2 every irreducible reduced fibre support
has nonnegative topological Euler number. For a
singular irreducible support, rational resolution
of the normalization gives strictly positive
Euler number; so does a resolution whose minimal
model is of class VII with positive second
Betti number. These assertions permit singular
normalization.

Every normal-crossing support has the
nonnegative even Euler number of P4.6.
Consequently, for a fibration all of whose
reduced supports are irreducible or have
normal crossings,
\[
 B=\mathbb P^1,\qquad c_3(Y)\in\{0,2\},
 \qquad b_3(Y)=2-c_3(Y).
\]
This last conclusion does not assume a
classification of the equality cases for
irreducible nonnormal supports.
\end{proposition}

% Exact cited input: series:II / night:P4.32
\begin{proposition}[P4.32]\hypertarget{tag-02X5}{}\label{night:P4.32}\label{stmt:P4.32}


An irreducible support in Part~4.2 has Euler number zero exactly
in the following cases: it is smooth with a finite unramified
torus cover, or it is singular with smooth minimal elliptic ruled
normalization, or it is singular with smooth minimal Hopf
normalization. In the ruled double-horizontal case the conductor
image is a smooth elliptic curve; the local germs are product
cusps, product tacnodes, or product ordinary plane triple points,
and the reduced conductor maps have degrees one, two, or three
respectively and are unramified.

For a fibration whose supports are all irreducible or normal
crossing, $c_3(Y)=0$ exactly when every irreducible support has
one of these forms and every reducible support has the
zero-Euler normal-crossing form of P4.8.
\end{proposition}

% Exact cited input: series:II / night:P4.35
\begin{proposition}[P4.35]\hypertarget{tag-02X8}{}\label{night:P4.35}\label{stmt:P4.35}


No irreducible support with rationally resolved normalization
has Euler number one. Consequently, in a fibration whose
supports are all irreducible or normal crossing, every
irreducible singular support with rationally resolved
normalization has Euler number exactly two.

More generally, an irreducible support cannot have reduced
image conductor equal to a rational curve with exactly one
point-identification cycle when its reduced normalization
conductor has no graph cycles. This excludes that conductor
case for every ruled-resolution double-horizontal pole tree.
\end{proposition}

% Exact cited input: series:II / night:P4.37
\begin{proposition}[P4.37]\hypertarget{tag-02XA}{}\label{night:P4.37}\label{stmt:P4.37}


No irreducible singular support in Part~4.2 has normalization
whose minimal resolution has a class-VII minimal model with
positive second Betti number. Consequently every class-VII
normalization occurring here is a smooth minimal Hopf surface
and has Euler number zero.
\end{proposition}

% Exact cited input: series:II / night:P4.54
\begin{proposition}[P4.54]\hypertarget{tag-02XT}{}\label{night:P4.54}\label{stmt:P4.54}


In the remaining positive ruled support of Part~4.2,
the horizontal pole section is not contracted by the
normalization resolution. It consequently maps with
degree two or three onto a rational normalized conductor
component whose generic transverse germ is $D_4$.
No transverse-type restriction is imposed on the other
conductor component.
\end{proposition}

% Exact cited input: series:II / night:P4.57
\begin{proposition}[P4.57]\hypertarget{tag-02XW}{}\label{night:P4.57}\label{stmt:P4.57}


The remaining degree-two configuration of P4.55--P4.56
cannot have its image-conductor graph cycle realized at
two distinct intersections of its two components.
\end{proposition}

% Exact cited input: series:II / night:P4.58
\begin{proposition}[P4.58]\hypertarget{tag-02XX}{}\label{night:P4.58}\label{stmt:P4.58}


Let $f:Y\to B$ be the fibration of Part~4.2, with
$b_1(Y)=b_2(Y)=0$, connected fibres and general fibre a
complex two-torus. Suppose every reduced fibre support
is irreducible or has normal crossings. Then
\[
 B=\mathbb P^1,\qquad c_3(Y)\in\{0,2\},
                       \qquad b_3(Y)=2-c_3(Y).
\]
Every reduced support has nonnegative Euler number.
An irreducible singular support whose normalization has
a rational resolution has Euler number two. Every other
irreducible support has Euler number zero. Such a support
is either smooth with a finite unramified torus cover,
or singular with smooth minimal elliptic ruled
normalization, or singular with smooth minimal Hopf
normalization. In particular the class-VII normalization
alternative is realized only by a smooth minimal Hopf
surface.

For a normal-crossing support the Euler number equals
its number of triple points, and is even. Its zero case
is exactly P4.8: in the reducible case the normalization
is a disjoint union of smooth minimal Hopf or non-K\"ahler
elliptic surfaces, with smooth elliptic conductor curves
and finite unramified branch maps. The equality
$c_3(Y)=0$ holds exactly when every support has one of
the corresponding zero-Euler forms.
\end{proposition}

% Exact cited input: series:II / night:P4R.18
\begin{proposition}[P4R.18]\hypertarget{tag-02YB}{}\label{night:P4R.18}\label{stmt:P4R.18}


In the remaining positive-Euler ruled case of Part~4.2, a
noncontracted horizontal pole section maps with degree two or
three to a rational normalized conductor component whose generic
transverse germ is $D_4$. This assertion imposes no transverse
type or conductor-exponent restriction on the other image component.
\end{proposition}

% Exact cited input: series:II / night:P4R.20
\begin{proposition}[P4R.20]\hypertarget{tag-02YE}{}\label{night:P4R.20}\label{stmt:P4R.20}


In the remaining degree-two horizontal sector after P4.55,
suppose that the generic transverse curve along the other
image component $B$ has at least three normalization branches.
The unique image-graph cycle cannot be concentrated at one
point with three local critical-curve branches, namely the
smooth branch $A$ and two local branches of $B$.
\end{proposition}

